% 补充图 C：DesignWare APB SSI 代码规模（学术风格：内刻度、黑描边、标签不越轴）
\documentclass[border=8pt]{standalone}
\usepackage{xeCJK}
\setCJKmainfont{Noto Serif CJK SC}
\usepackage{tikz}
\usepackage{xcolor}
\definecolor{figblue}{HTML}{1D4ED8}
\definecolor{figorange}{HTML}{F97316}
\begin{document}
\begin{tikzpicture}[x=0.0047cm,y=0.72cm,font=\footnotesize]
  % 轴与刻度（向内）
  \draw[line width=0.8pt] (0,0.25) -- (0,-5.7) -- (2120,-5.7);  \foreach \x in {0,500,1000,1500,2000} {
    \draw[line width=0.5pt] (\x,-5.7) -- (\x,-5.56);
    \node[below, font=\scriptsize] at (\x,-5.72) {\x};
  }
  \node[below=11pt, font=\scriptsize] at (1060,-5.72) {代码行数};

  % 横条（黑描边；值标签置于柱端外侧，不越轴）
  \path[fill=figblue, draw=black, line width=0.5pt]
    (0,-0.42) rectangle (1844,0.0);
  \node[right, font=\scriptsize] at (1848,-0.21) {1{,}844};
  \node[left, align=right, font=\footnotesize] at (-14,-0.21)
    {Linux 源码（输入，固定）\\ {\scriptsize 核心 + MMIO 粘合 + 头文件}};

  \path[fill=figorange, draw=black, line width=0.5pt]
    (0,-1.5) rectangle (2092,-1.08);
  \node[right, font=\scriptsize] at (2096,-1.29) {2{,}092};
  \node[left, align=right, font=\footnotesize] at (-14,-1.29)
    {生成 Linux C\\ {\scriptsize 源文件 + 头文件}};

  \path[fill=figorange, draw=black, line width=0.5pt]
    (0,-2.44) rectangle (1977,-2.02);
  \node[right, font=\scriptsize] at (1981,-2.23) {1{,}977};
  \node[left, align=right, font=\footnotesize] at (-14,-2.23)
    {生成 harness C\\ {\scriptsize 源文件 + 头文件}};

  \path[fill=figorange, draw=black, line width=0.5pt]
    (0,-3.38) rectangle (1946,-2.96);
  \node[right, font=\scriptsize] at (1950,-3.17) {1{,}946};
  \node[left, align=right, font=\footnotesize] at (-14,-3.17)
    {生成裸机 C\\ {\scriptsize 源文件 + 头文件}};

  \path[fill=figorange, draw=black, line width=0.5pt]
    (0,-4.32) rectangle (1909,-3.9);
  \node[right, font=\scriptsize] at (1913,-4.11) {1{,}909};
  \node[left, align=right, font=\footnotesize] at (-14,-4.11)
    {生成 Rust \texttt{no\_std}\\ {\scriptsize tock-registers}};

  % 注记（白底细框）
  \node[draw=black, line width=0.5pt, fill=white, align=left, font=\scriptsize,
    anchor=north west, inner sep=0.16cm, text width=10.6cm] at (-60,-6.85)
  {各目标行数与源码相当：每个目标约 316 行为回执/锚点标记（无功能内容），harness 对约占 316 行。
   目标只实现提取出的设备核心 + 该后端适配器，缩减的是 LLM 的语义与文本搜索空间，
   不是删改源码，也不构成正确性论证（§1、§6.6）。};
\end{tikzpicture}
\end{document}
